\documentclass[11pt]{article}

\usepackage[a4paper,margin=25mm]{geometry}
\usepackage{amsmath,amssymb,amsthm}
\usepackage{graphicx}
\usepackage{microtype}
\usepackage[hidelinks]{hyperref}

\newtheorem{theorem}{Theorem}
\newtheorem{proposition}{Proposition}
\newtheorem{corollary}{Corollary}
\theoremstyle{definition}
\newtheorem{definition}{Definition}
\theoremstyle{remark}
\newtheorem*{remark}{Remark}

\newcommand{\C}{\mathbb C}
\newcommand{\D}{\mathbb D}
\newcommand{\Nzero}{\mathbb N_0}
\newcommand{\Hcal}{\mathcal H}
\newcommand{\Mult}{\operatorname{Mult}}
\newcommand{\norm}[1]{\left\lVert #1\right\rVert}

\title{A Superexponentially Weighted Hilbert Space\\
with Only Constant Multipliers}
\author{Candidate solution to Question 4.9 of arXiv:2008.03700}
\date{9 August 2026}

\begin{document}
\maketitle

\begin{abstract}
We give an affirmative answer to Bilokopytov's Question 4.9.  The Hilbert
space of power series with coefficient norm
\[
  \norm{f}_{\Hcal}^{2}=\sum_{n\geq0}|a_n|^2 e^{2n^2},
  \qquad f(z)=\sum_{n\geq0}a_nz^n,
\]
consists in fact of entire functions, so it is a Hilbert space of holomorphic
functions over the unit disk.  If a multiplier has a nonzero Taylor
coefficient of degree $j\geq1$, its action on normalized monomials has norm at
least $|b_j|e^{2jn+j^2}$, which is unbounded as $n\to\infty$.  Consequently
the multiplier algebra is exactly $\C$.
\end{abstract}

\section*{Source Problem}

The source is Eugene Bilokopytov, \emph{Multiplier algebras of normed spaces
of continuous functions}, arXiv:2008.03700 \cite{Bilokopytov2020}.  The open
problem occurs as Question 4.9 on page 13 of the arXiv PDF, immediately after
Corollary 4.8.

\begin{center}
  \includegraphics[width=\linewidth]{figures/open_problem_crop.png}
\end{center}

The question reads:
\begin{quote}
Does there exist an anti-multiplicative HSHF over $\D$?
\end{quote}

\section*{Definitions}

Following the source, a \emph{Hilbert space of holomorphic functions} (HSHF)
over $\D$ is a Hilbert space $\Hcal\subseteq\mathcal O(\D)$ whose norm
topology is stronger than the compact-open topology.  Equivalently in the
Hilbert setting, all point evaluations are bounded; the closed graph/uniform
boundedness argument then gives the required compact-open control.

A function $\varphi:\D\to\C$ is a \emph{multiplier} of $\Hcal$ if
$M_\varphi f=\varphi f$ defines a bounded operator $M_\varphi:\Hcal\to\Hcal$.
The space is \emph{anti-multiplicative} when
$\Mult(\Hcal)=\C\mathbf 1$, i.e. its only multipliers are constants.

\section*{Proof Intuition}

In a coefficient Hilbert space with monomial weights $w_n$, multiplication by
$z^j$ shifts degree $n$ to degree $n+j$ and has size $w_{n+j}/w_n$ on the
normalized $n$th monomial.  We choose $w_n=e^{n^2}$.  Then every fixed
positive shift has exploding ratio
\[
  \frac{w_{n+j}}{w_n}=e^{(n+j)^2-n^2}=e^{2jn+j^2}\longrightarrow\infty.
\]
For an arbitrary nonconstant multiplier, isolate any one nonzero positive
Taylor coefficient.  Its contribution to a single output coefficient already
forces exactly the same explosion; no cancellation can occur because each
input is a monomial.

The only other obligation is to make the Hilbert space genuinely a function
space on $\D$.  The reciprocal weights decay faster than every geometric
sequence, so Cauchy--Schwarz gives uniform convergence on every compact subset
of the whole plane.

\section*{Main Result}

\begin{theorem}\label{thm:main}
There exists a $1$-containing, infinite-dimensional HSHF $\Hcal$ over $\D$
such that
\[
  \Mult(\Hcal)=\C\mathbf 1.
\]
In particular, Question 4.9 of \cite{Bilokopytov2020} has an affirmative
answer.
\end{theorem}

\begin{proof}
Define
\[
  \Hcal=\left\{f(z)=\sum_{n=0}^{\infty}a_nz^n:
  \norm{f}_{\Hcal}^{2}:=\sum_{n=0}^{\infty}|a_n|^2e^{2n^2}<\infty\right\}.
\]
The coefficient map
\[
  f\longmapsto (a_ne^{n^2})_{n\geq0}
\]
is an isometry onto $\ell_2(\Nzero)$, so $\Hcal$ is a Hilbert space.  To see
that its elements define holomorphic functions with the required topology,
fix $R>0$.  Cauchy--Schwarz gives, for $|z|\leq R$,
\begin{align*}
 |f(z)|
 &\leq \sum_{n\geq0}|a_n|R^n \\
 &\leq \left(\sum_{n\geq0}|a_n|^2e^{2n^2}\right)^{1/2}
       \left(\sum_{n\geq0}R^{2n}e^{-2n^2}\right)^{1/2} \\
 &= C_R\norm{f}_{\Hcal},
\end{align*}
where $C_R<\infty$ by the ratio test.  Hence the defining power series
converges uniformly on every disk $|z|\leq R$; every element of $\Hcal$ is
entire.  The displayed estimate also makes the inclusion into the
compact-open topology continuous (in particular on compact subsets of
$\D$).  Thus $\Hcal$ is an HSHF over $\D$.  It contains $\mathbf1$, and it is
infinite-dimensional because it contains all monomials.

Let $\varphi$ be a multiplier of $\Hcal$.  Since $\mathbf1\in\Hcal$,
\[
  \varphi=M_\varphi\mathbf1\in\Hcal.
\]
Write its Taylor expansion as
\[
  \varphi(z)=\sum_{k=0}^{\infty}b_kz^k.
\]
Suppose that $\varphi$ is nonconstant.  Choose $j\geq1$ with $b_j\neq0$.
For each $n\geq0$, put
\[
  u_n(z)=e^{-n^2}z^n.
\]
Then $\norm{u_n}_{\Hcal}=1$.  The coefficient of $z^{n+j}$ in
$\varphi u_n$ is exactly $b_je^{-n^2}$.  Since the squared Hilbert norm is a
sum of nonnegative coefficient contributions,
\begin{align*}
  \norm{M_\varphi u_n}_{\Hcal}
  &\geq |b_j|e^{-n^2}e^{(n+j)^2} \\
  &=|b_j|e^{2jn+j^2}\longrightarrow\infty.
\end{align*}
This contradicts the boundedness of $M_\varphi$ on the unit vectors $u_n$.
Therefore every multiplier is constant.  Conversely, each constant $c$ acts
as the bounded operator $cI_{\Hcal}$.  Hence
$\Mult(\Hcal)=\C\mathbf1$.
\end{proof}

\begin{remark}[Reproducing kernel]
The construction may equivalently be presented as the reproducing-kernel
Hilbert space with kernel
\[
  K(z,w)=\sum_{n=0}^{\infty}e^{-2n^2}(z\overline w)^n.
\]
The series converges for all $(z,w)\in\C^2$, and its associated orthonormal
basis is $e^{-n^2}z^n$.  This makes explicit that the example is an ordinary
RKHS, not merely a formal coefficient space.
\end{remark}

\section*{Verification Notes}

The proof has four independent checks.
\begin{enumerate}
  \item Completeness is exactly completeness of $\ell_2$ under the displayed
        coefficient isometry.
  \item The Cauchy--Schwarz estimate holds for every $R>0$, a statement
        stronger than needed for an HSHF over $\D$.
  \item The test vectors are normalized: their only coefficient has weighted
        magnitude $e^{-n^2}e^{n^2}=1$.
  \item There is no cancellation in the lower bound.  For a fixed monomial
        input, the coefficient of output degree $n+j$ comes from the unique
        Taylor coefficient $b_j$; discarding all other nonnegative squared
        coefficient contributions gives a valid lower bound.
\end{enumerate}
No numerical or computer-assisted claim is used.

\section*{Novelty Check And Limitations}

Before promotion, the run's registry and solution, attempt, and proof-gap
indexes were searched for arXiv:2008.03700, its title, and the phrases
``anti-multiplicative HSHF'', ``Hilbert space of holomorphic functions'', and
``only constant multipliers''.  No duplicate result or prior attempt was
found.

A bounded web search on 9 August 2026 used the exact text of Question 4.9,
the exact term ``anti-multiplicative HSHF'', and close variants involving
reproducing-kernel Hilbert spaces with only constant multipliers.  The search
located the source paper and general background results, but no separate later
paper explicitly answering Question 4.9.  This was not an exhaustive
bibliographic or citation-graph search, so novelty remains subject to human
literature review.

The result answers the existence question exactly as stated.  It does not
claim additional regularity of the norm under classical operators such as the
shift, nor does it classify all HSHFs with trivial multiplier algebra.  The
rapid weights are the mechanism that deliberately make every positive-degree
shift unbounded.

\section*{Human Review Recommendation}

Promote for expert review as a candidate full affirmative answer.  The proof
is short and self-contained.  A reviewer should verify the source convention
for HSHF and the coefficient-isolation argument for arbitrary multiplier
symbols.  Under the definitions in \cite{Bilokopytov2020}, both checks are
direct: the compact-open embedding is explicit, $\mathbf1\in\Hcal$, and
Taylor coefficients are unique.

\begin{thebibliography}{9}
\bibitem{Bilokopytov2020}
E.~Bilokopytov,
\emph{Multiplier algebras of normed spaces of continuous functions},
arXiv:2008.03700, 2020.
\url{https://arxiv.org/abs/2008.03700}.
\end{thebibliography}

\end{document}
